% report_facct.tex — FAccT/ACM first draft. PRIMARY paper for this project.
%
% The NeurIPS Datasets & Benchmarks variant is report/report.tex. The two make
% different claims and are not two formats of one paper: report.tex argues that
% the benchmark is a reusable artifact; this paper argues that fairness audit
% infrastructure encodes normative commitments which surface as penalties on the
% correct answer, and uses the benchmark as method.
%
% Build:  make paper-facct   (tectonic; it runs BibTeX itself)
%
% For double-blind submission, switch the documentclass line to:
%   \documentclass[sigconf,anonymous,review]{acmart}
% and remove \author blocks' affiliation detail. `nonacm` is set because this is
% a draft with no assigned DOI/ISBN; drop it once the venue supplies them.

\documentclass[sigconf,nonacm]{acmart}

\usepackage{booktabs}
\usepackage{multirow}

\citestyle{acmnumeric}

\begin{document}

\title{Penalising the Right Answer}
\subtitle{How Fairness Audit Infrastructure Encodes Normative Commitments,
and What Three Lending Audits Revealed}

\author{Siddhant Jain}
\affiliation{%
  \institution{Simon Fraser University}
  \city{Burnaby}
  \state{British Columbia}
  \country{Canada}}
\email{sja178@sfu.ca}

\author{Ayaan Lalani}
\affiliation{%
  \institution{Simon Fraser University}
  \city{Burnaby}
  \state{British Columbia}
  \country{Canada}}

\begin{abstract}
A fairness audit is a claim about a decision that affects people. The
infrastructure that produces such claims --- metric libraries, reporting floors,
scoring harnesses, quality checklists --- is usually treated as neutral
apparatus. It is not: it infers normative categories from syntactic features,
and the inference fails hardest on the inputs where judgement was most needed.
We demonstrate this four times over a lending audit, twice against instruments we
built ourselves.

A subgroup reporting floor of $n \geq 15$, a standard and defensible statistical
commitment, suppressed five of eleven cells in a mortgage audit. Three of them
were worse than anything the audit reported, so it named the wrong group as most
disadvantaged and understated severity by two bands: it returned
\textsc{moderate} at $0.8169$ while the worst real disparity was $0.1917$.
Auditing the full population instead of a test split makes those groups visible
and moves the worst reported cell to $0.6734$. A refusal detector flagged six
correct answers as refusals because they fell below a character threshold under a
prompt demanding terseness; applied to a second provider it mechanically zeroed
fourteen cycles, moving that model's mean score from $70.0$ to $19.79$, and a
permanent block was written into our configuration on that evidence. A
fabricated-metric detector flagged a newer model five times as often as its
predecessor --- but $67\%$ of those flags are values derivable from the provided
context, mostly ratios restated as percentages, so the better-reasoning model is
penalised for showing its work. And auditing age on the $40+$ threshold borrowed
from employment law returns near parity ($0.9791$) where the cut the credit
statute specifies, $62+$, returns the largest age disparity in the data
($0.8897$).

We also report what this cost us. Three previously published findings of ours do
not survive: a $\textsc{critical}$ intersectional ratio of $0.6087$ that becomes
$0.8201$ and clears the four-fifths screen once its fifteen-person comparator
grows; an inference about label bias drawn from a classifier with target leakage
(AUC $0.9942 \to 0.7915$ once repaired); and an arithmetic claim that $0.8956$
failed a $0.80$ screen. We report the retractions because the retracted versions
were the more quotable ones.
\end{abstract}

\begin{CCSXML}
<ccs2012>
<concept>
<concept_id>10003456.10003462.10003480</concept_id>
<concept_desc>Social and professional topics~Governmental regulations</concept_desc>
<concept_significance>300</concept_significance>
</concept>
<concept>
<concept_id>10010147.10010257</concept_id>
<concept_desc>Computing methodologies~Machine learning</concept_desc>
<concept_significance>500</concept_significance>
</concept>
<concept>
<concept_id>10003120.10003121</concept_id>
<concept_desc>Human-centered computing~Human computer interaction (HCI)</concept_desc>
<concept_significance>300</concept_significance>
</concept>
</ccs2012>
\end{CCSXML}

\ccsdesc[500]{Computing methodologies~Machine learning}
\ccsdesc[300]{Social and professional topics~Governmental regulations}
\ccsdesc[300]{Human-centered computing~Human computer interaction (HCI)}

\keywords{algorithmic fairness, audit, disparate impact, intersectionality,
lending, large language models, measurement}

\maketitle

\section{Introduction}

A disparate impact ratio of $0.86$ is not a finding. It becomes one when attached
to who was denied credit, what recourse they had, and what the lender is obliged
to check. \citet{lee2020} make this the organising commitment of human-centered
fair AI: fairness is defined relative to a human decision context, not in metric
space. We adopt that commitment and follow it one step further than is usual. If
a fairness claim depends on its decision context, then so does the apparatus that
produces the claim --- and that apparatus can be wrong in ways that are invisible
from inside it.

This is not a hypothetical. \citet{deng2022} find practitioners using fairness
toolkits in ways their designers did not anticipate, and \citet{leesingh2021} map
systematic gaps in what open-source toolkits actually support.
\citet{holstein2019} report that practitioners' needs diverge sharply from what
the research literature supplies. \citet{mei2025} show, in a speech setting, that
audit \emph{pipelines} introduce pitfalls that change conclusions.
\citet{zaccour2025} show that quantitative audits are more often blocked by data
access than by method. The common thread is that the instrument is part of the
finding.

We built a fairness audit across three lending decisions --- consumer credit
approval, mortgage origination, and peer-to-peer default risk --- with a
deterministic metric layer and a language model supplying the interpretive
narrative. The intent was a benchmark. What the exercise produced, and what we
report here, is a set of findings about audit infrastructure, three of which
implicate our own instruments.

\paragraph{Claim.} Fairness audit infrastructure infers normative categories from
syntactic features --- a label value, a character count, a row count, a threshold
borrowed from another statute --- and the inference fails hardest on the inputs
where judgement was most needed. We support this with four demonstrations:

\begin{enumerate}
  \item \textbf{Metric orientation} (\S\ref{sec:orientation}). Correct toolkit,
        clean data, sound metric, and a maximally wrong conclusion --- disparate
        impact $0.0$ where the true value is $0.9931$ --- produced by one
        convention about which label counts as favorable.
  \item \textbf{Reporting floors} (\S\ref{sec:floor}). An $n \geq 15$ floor
        suppressed five of eleven cells, three of them worse than anything the
        audit reported. The cost of a prudent statistical threshold falls
        entirely on the smallest groups.
  \item \textbf{The wrong statute} (\S\ref{sec:ecoa}). Auditing age at $40+$, a
        threshold from employment law, returns a clean result on the one age
        class credit law actually protects.
  \item \textbf{Mechanical proxies for judgement} (\S\ref{sec:instruments}).
        Three instruments we wrote in good faith systematically punished the
        correct response. We did not repair them, because repairing them would
        have hidden the finding.
\end{enumerate}

\noindent We further report (\S\ref{sec:interpretive}) that the interpretive
layer is contestable, and that its instability is model-generation-specific:
prompt phrasing moves severity verdicts on \texttt{gpt-4o} by up to $14.7$ points
and by $0.0$ on \texttt{gpt-5-mini}, on identical input.

\paragraph{What we retract.} Three findings we previously published do not
survive this re-analysis, and we set them out in place rather than quietly
dropping them (\S\ref{sec:floor}, \S\ref{sec:intersectional}). Audit
infrastructure that cannot be corrected in public is the problem this paper is
about.

\paragraph{Position.} We are not arguing that fairness metrics are useless or
that audit tooling should be distrusted wholesale. We are arguing something
narrower and more actionable: the failure modes documented here are not user
error, and they are not exotic. Each arises from a defensible design decision
made by someone who understood the domain. That is what makes them worth
reporting at this venue.

\section{Related work}

\paragraph{Human-centered framing.} \citet{lee2020} is our framing exemplar: a
fairness claim is meaningful only relative to a human decision context. The line
of work extending it has moved steadily toward stakeholder involvement in the
choice of fairness criteria rather than in the interpretation of a fixed metric.
\citet{luo2025} develop a protocol for explaining, negotiating and agreeing
fairness metrics with stakeholders; \citet{luo2026} document how much
disagreement stakeholders exhibit even when shown identical evidence;
\citet{luo2026b} extend this to affected individuals rather than institutional
stakeholders. \citet{deshpande2022} ask who the stakeholders of a responsible AI
system even are. \citet{stumpf2024} argue that fairness and correctness need
user-centred assessment rather than metric reporting. \citet{kim2025} survey how
equity is operationalised across HCI and fairness research and find the construct
fragmenting under practical demand.

Our contribution is adjacent rather than competing. This literature asks who
should choose the criteria. We ask what the apparatus does once a criterion has
been chosen, and show that the apparatus can invert the answer without anyone
choosing anything.

\paragraph{Audit practice and its instruments.} \citet{metaxa2021} survey
algorithm auditing as a method. \citet{raji2019} demonstrate that external audits
change vendor behaviour, and \citet{raji2020} give an internal end-to-end audit
framework. \citet{kroll2021} argues for traceability as the operational principle
behind accountability, which is the design stance our frozen-artifact approach
takes. Documentation interventions --- model cards \citep{mitchell2019},
datasheets \citep{gebru2021} --- and process interventions such as co-designed
checklists \citep{madaio2020} share our premise that fairness work needs
infrastructure, and share the risk we report: a checklist is a mechanical proxy
for judgement, and \S\ref{sec:instruments} is a case study in what that costs.

\citet{mei2025} is the nearest prior work: audit pitfalls in an ASR setting that
change the conclusion. \citet{zamfirescu2022} diagnose human error inside an
image-analysis pipeline. We add a lending case in which the pitfall is a metric
convention and the affected quantity is the headline number.

\paragraph{Measurement.} That fairness criteria conflict is settled
\citep{chouldechova2017,kleinberg2017}, and that the choice among them is not a
technical matter is argued forcefully by \citet{selbst2019}, whose abstraction
traps describe our setting well. \citet{corbettdavies2018} catalogue how fairness
measures mislead when detached from decision context. \citet{barocas2016} and
\citet{wachter2021} establish that the legal notion of discrimination does not
reduce to a metric threshold --- which is why we treat the four-fifths screen as
a screening signal throughout and make no legal claim. On intersectionality,
\citet{buolamwini2018} provide the canonical empirical demonstration that
marginal accuracy hides joint-group failure, and \citet{foulds2020} give a formal
definition. \S\ref{sec:intersectional} contributes a lending instance where the
marginal audit passes cleanly and the joint audit does not.

\paragraph{Language models as auditors and judges.} Using an LLM to interpret
metrics invites the LLM-as-judge critique \citep{zheng2023}: model-graded
evaluation inherits the grader's biases. Track Q avoids the worst of this: its
ground truth is a deterministic rule over computed metrics, never another
model's opinion (\S\ref{sec:method}). The one cross-model grade we report is
flagged as secondary evidence (\S\ref{sec:interpretive}). Most directly,
\citet{wester2024} study when and how LLMs deny user requests, which is the exact
phenomenon our refusal detector was built to catch and, as
\S\ref{sec:instruments} reports, systematically failed to measure.

\section{Setting and method}
\label{sec:method}

The benchmark is method here, not contribution; the artifact claim is made
elsewhere. What matters for this paper is that every number an LLM saw was
computed by deterministic code, and that the model was never permitted to
recompute one.

\subsection{Three lending decisions}

The three use cases are deliberately not three instances of one setting
(Table~\ref{tab:usecases}). They differ in the two properties that most change
what an audit can conclude: whether the protected attribute is a protected class
or an economic proxy, and whether the classifier discriminates enough for
disparity to be measurable at all. Lending discrimination is documented
empirically across this span --- appearance and gender effects in peer-to-peer
markets \citep{duarte2012,pope2011}, FinTech-era consumer lending
\citep{bartlett2019}, local credit-market conditions \citep{butler2017}, and
historical records used as training data \citep{sarwal2024}.

\begin{table}[t]
\caption{The three use cases. The last row governs how every downstream result
reads.}
\label{tab:usecases}
\small
\begin{tabular}{@{}p{1.35cm}p{1.75cm}p{1.75cm}p{1.9cm}@{}}
\toprule
& \textbf{UC1 German Credit} & \textbf{UC2 HMDA (GA)} & \textbf{UC3 Lending Club}\\
\midrule
Decision & consumer credit & mortgage origination & P2P default risk\\
Attributes & sex, age, foreign worker & race, sex, age & gender, income, loan size\\
Type & protected class & protected class & 1 class + 2 proxies\\
Favorable & good credit (1) & approved (1) & non-default (\textbf{0})\\
Classifier & discriminating & discriminating & \textbf{$\sim$99.5\% one class}\\
\bottomrule
\end{tabular}
\end{table}

\subsection{Deterministic layer}

For each protected attribute we compute disparate impact (DI), demographic parity
difference (DPD), equal-opportunity difference (EOD), average odds difference
(AOD) and the Theil index with AIF360 \citep{aif360}, cross-checked against
Fairlearn \citep{fairlearn}. A fixed rule maps metric values to a severity label
(\textsc{low}, \textsc{moderate}, \textsc{high}, \textsc{critical}); the
\textsc{critical} boundary for DI is $0.72$. Subgroup cells below $n = 15$ are
reported with their $n$ and excluded from conclusions. Data are German Credit
\citep{germancredit}, HMDA Georgia \citep{hmda}, and a Lending Club
peer-to-peer default task.

Two conventions are pinned by regression tests because both bit us: which label
counts as favorable (\S\ref{sec:orientation}), and the sign of AOD.

\subsection{Two tracks, kept separate}

The interpretive layer is evaluated along two axes that we never score against
each other, because they answer different questions.

\textbf{Track Q} asks whether the model interprets fixed metrics reliably. For
each of nine (use case $\times$ attribute) pairs we issue four prompts ---
\texttt{zero\_shot}, \texttt{chain\_of\_thought}, \texttt{self\_critique},
\texttt{constrained} --- yielding 36 \texttt{gpt-4o} calls, each scored against
the deterministic severity label. All 36 prompts are frozen and released, and a
replay path reproduces the entire scoring pipeline with no API key.

\textbf{Track H} asks whether the narrative form raises the ceiling on audit
quality. Three \texttt{o3} audits, one per use case, judged only against a
five-criterion checklist, a guardrail baseline, and a quality rubric --- never
against the deterministic severity numbers, which would merely re-run Track Q.

Two detectors run over every response: a fabricated-citation check (n-gram match
against the titles actually retrieved) and a fabricated-metric check (any numeric
claim absent from the provided context). Spend is capped cumulatively at \$15 with a
pre-call abort and a persisted ledger. The ledger records \textbf{98 calls
totalling \$0.8933}; an additional \$0.1125 was spent on a pre-freeze run whose
calls predate the ledger. We state calls and dollars separately because an
earlier version of this paper conflated them.

\section{Finding 1: metric orientation inverts the headline number}
\label{sec:orientation}

In UC3, \texttt{gender} initially returned $\mathrm{DI} = 0.0$ and a
\textsc{critical} severity label --- the most severe output the pipeline can
produce, and on its face evidence of near-total exclusion.

It was an artifact. The prediction target is \texttt{loan\_default}, so label $1$
is the \emph{unfavorable} outcome. Computing selection rates on predicted default
gives $0.0000$ for men and $0.0069$ for women, and
$\mathrm{DI} = 0.0000/0.0069 = 0.0$. Oriented on the favorable outcome ---
predicted non-default --- the rates are $1.0000$ and $0.9931$, so
$\mathrm{DI} = 0.9931$: near parity (Table~\ref{tab:uc3}).

\begin{table}[t]
\caption{UC3 Lending Club, oriented on the favorable outcome (predicted
non-default, label $0$). \texttt{gender} previously read $\mathrm{DI}=0.0$,
\textsc{critical}.}
\label{tab:uc3}
\small
\begin{tabular}{@{}lrrrr@{}}
\toprule
Attribute & DI & DPD & EOD & AOD\\
\midrule
\texttt{gender} (priv.\ male) & \textbf{0.9931} & $-0.0069$ & $0.0$ & $-0.0357$\\
\texttt{income\_level} (priv.\ med.) & 1.0072 & 0.0072 & $0.0$ & 0.0323\\
\texttt{loan\_amount\_level} (priv.\ med.) & 1.0063 & 0.0062 & $0.0$ & 0.0270\\
\bottomrule
\end{tabular}
\end{table}

Three things make this worth a section rather than a footnote.

\textbf{Nothing was misused.} AIF360 computes exactly what it is asked to. The
data were clean. Disparate impact is the appropriate screen for this question.
The error lives entirely in the mapping from a column name to a normative
category, and that mapping is a convention rather than a computation. A target
named \texttt{approved} makes label $1$ favorable; a target named
\texttt{loan\_default} inverts it. Nothing in the metric signature records which
situation holds.

\textbf{The failure is maximally severe and silent.} The reported value was wrong
by a factor limited only by floating point, and it was wrong in the alarming
direction, which is the direction least likely to prompt a second look: a
\textsc{critical} disparity on a protected attribute reads as the audit working.
Had the error run the other way it would have produced a false clean bill of
health, which is the same defect with worse consequences.

\textbf{It generalises.} Default, delinquency, charge-off, fraud and churn are
all conventionally coded with the adverse event as $1$. Any fairness audit of a
risk model --- which is to say most fairness audits in lending --- meets this
condition. \citet{deng2022} and \citet{leesingh2021} document toolkit-practice
gaps of exactly this kind, and \citet{selbst2019} name the general form: the
formalism omits the part of the problem that determines the answer.

\paragraph{Implication.} Favorable-outcome orientation should be an explicit,
required argument in fairness tooling, not inferred from label values or left to
the analyst. It is now pinned by regression tests in our pipeline, which converts
a silent inversion into a build failure. That is a small change with a large
blast radius, and it is available to any toolkit maintainer today.

\section{Finding 2: the reporting floor erases the largest disparities}
\label{sec:floor}

Subgroup reporting floors are standard practice and defensible: a ratio computed
on six people is not evidence. Ours suppresses cells below $n = 15$. What that
commitment does, in our data, is remove the most disadvantaged groups from the
conclusion.

We hold the model fixed and vary only the audited population --- a held-out test
split, versus every row of the cleaned dataset scored out-of-fold by a model that
never saw it. This isolates sample size from every other change
(Table~\ref{tab:floor}).

\begin{table}[t]
\caption{HMDA \texttt{race} $\times$ \texttt{sex} under an $n \geq 15$ floor.
Identical leak-free model; only the audited population differs.}
\label{tab:floor}
\small
\begin{tabular}{@{}lrr@{}}
\toprule
& Test split & Full population\\
& ($n{=}2{,}196$) & ($n{=}10{,}978$)\\
\midrule
Cells suppressed & \textbf{5 of 11} & 1 of 12\\
Worst cell \emph{reported} & 0.8169 (\textsc{mod.}) & 0.6734 (\textsc{crit.})\\
Worst cell that \emph{exists} & \textbf{0.1917} (\textsc{crit.}) & 0.6734 (\textsc{crit.})\\
Suppressed cells worse & \textbf{3} & 0\\
\quad than any reported & & \\
\bottomrule
\end{tabular}
\end{table}

On the test split the audit returns \textsc{moderate} for
\texttt{Black $\times$ Male} at $0.8169$ while \texttt{American Indian $\times$
Female} sits at $0.1917$, unreported. The audit is wrong about which group is
most disadvantaged and wrong by two severity bands about how bad it is. Every
suppressed cell in that run belongs to a racial minority.

At full population the same floor suppresses one cell of twelve, and the worst
disparity --- \texttt{American Indian $\times$ Male}, $0.6734$, $n=31$ ---
becomes visible. Nothing about the population changed; the audit simply stopped
being run on a sample too small to see it.

The floor is not a bug, and we do not propose removing it. The point is narrower
and harder to design around: a threshold chosen for statistical prudence
distributes its cost entirely onto the smallest groups, which in a lending
context are the groups least able to absorb it. Two further details make it worse
in practice. Our own codebase carries \emph{two} floors, $n \geq 15$ in the
drilldown and $n \geq 30$ in the qualitative layer, so the same cell is reportable
or not depending on which module reads it. And the floor is applied silently: the
audit output lists what passed, not what was withheld.

\subsection{The marginal-versus-intersectional gap, corrected}
\label{sec:intersectional}

Every marginal attribute in German Credit clears the four-fifths screen
(Table~\ref{tab:uc1}). An audit reporting per-attribute disparate impact would
return three passes and stop.

\begin{table}[t]
\caption{UC1 German Credit marginals. All three clear the four-fifths screen.}
\label{tab:uc1}
\small
\begin{tabular}{@{}lrrrr@{}}
\toprule
Attribute & DI & DPD & EOD & AOD\\
\midrule
\texttt{Sex} (priv.\ male) & 0.8888 & --- & --- & ---\\
\texttt{AgeGroup} (priv.\ 40+) & 0.9258 & --- & --- & ---\\
\texttt{foreign\_worker} (priv.\ 0) & 0.8387 & --- & --- & ---\\
\bottomrule
\end{tabular}
\end{table}

The intersection is the worst cell: \texttt{under\_40 $\times$ female} has
$\mathrm{DI} = 0.8201$ ($0.7012$ against $0.8551$ for
\texttt{40\_plus $\times$ female}), lower than any marginal attribute.

\paragraph{What we previously reported here, and why it was wrong.} On the
200-row test split this cell read $\mathrm{DI} = 0.6087$, which we published as
\textsc{critical}. Its comparator was \texttt{40\_plus $\times$ female} at
$n = 15$ with a selection rate of exactly $1.0000$ --- fifteen people, all
approved. A ratio against a saturated fifteen-person cell is not a severity
grade. At full population that comparator holds 69 people at $0.8551$ and the
ratio rises to $0.8201$, which \emph{clears} the four-fifths screen.

The direction survives and the magnitude does not. We had also written that the
disadvantaged group ``is not identified by either single-attribute audit'';
that was too strong, since \texttt{under\_40 $\times$ male} at $0.7952$ failed
the screen on the same split, so the age effect was visible marginally. The
honest claim is that marginal auditing understated a disparity it did point at.

This is the same defect as \S\ref{sec:floor} seen from the other side: there, a
small cell was suppressed and the disparity vanished; here, a small cell was
retained as a \emph{reference} and the disparity inflated. A floor governs which
cells are reported. Nothing in the pipeline governs which cell becomes the
comparator, and that choice moved a published verdict across two severity bands.

\subsection{UC2: the joint disparity exceeds both marginals}

HMDA shows the same structure (Table~\ref{tab:uc2}). At full population
\texttt{Black $\times$ Female} ($n = 1{,}994$) sits at $\mathrm{DI} = 0.8632$,
below both contributing marginals --- the joint \emph{disparity} exceeds the
marginal ones, though the ratio is numerically lower, a phrasing we previously
inverted.

Two caveats we did not previously state. First, the reference cell is chosen as
the most-favoured adequately-sized subgroup, which is a winner's-curse
estimator: measured instead against \texttt{White $\times$ Male}, the control
group conventional in US fair lending analysis, the same cell reads $0.8388$ on
the test split and clears the screen. Reference-group choice alone moves the
verdict across the four-fifths line, and nothing in the metric signature records
which convention is in force. Second, and following \S\ref{sec:floor}, this cell
is not the worst one; it is merely the worst \emph{large} one.

\begin{table}[t]
\caption{UC2 HMDA (Georgia) marginals, leak-free model, full population
($n{=}10{,}978$). The banded \texttt{age\_group} reads near parity while
\texttt{age\_62\_plus} --- the cut ECOA actually specifies --- is the lowest of
the four.}
\label{tab:uc2}
\small
\begin{tabular}{@{}lrrrr@{}}
\toprule
Attribute & DI & DPD & EOD & AOD\\
\midrule
\texttt{race} (priv.\ White) & 0.9371 & --- & --- & ---\\
\texttt{sex} (priv.\ Male) & 0.9572 & --- & --- & ---\\
\texttt{age\_group} (priv.\ mid, banded) & 0.9791 & --- & --- & ---\\
\texttt{age\_62\_plus} (priv.\ under 62) & \textbf{0.8897} & --- & --- & ---\\
\bottomrule
\end{tabular}
\end{table}

\subsection{Two readings a threshold gate destroys}

The same tables carry two conclusions that survive only if the metrics are read
together rather than screened individually.

\textbf{A withdrawn inference.} We previously read UC2's near-zero
$\mathrm{EOD}$ and $\mathrm{AOD}$ alongside a low $\mathrm{DI}$ as evidence of
label bias rather than differential model error. That reading is withdrawn. It
rested on a classifier trained with \texttt{interest\_rate} among its features, a
column populated for 100.0\% of originations and 0.7\% of denials: the model was
substantially detecting whether an interest rate existed, which is the outcome.
Its errors were near zero in every group because it made almost no errors at
all. Removing the leak drops AUC from $0.9942$ to $0.7915$. We report the
retraction rather than the inference, because the inference was the more
quotable of the two.

\textbf{A non-finding is not a pass.} UC1 \texttt{foreign\_worker} has
$\mathrm{DI} = 0.9402$, comfortably clear. The privileged group has \emph{six
records}, against $n = 194$ on the other side and a reporting floor of
$n \geq 15$. The ratio is dominated by the larger group and is not evidence of
anything. Reported as a pass, it is worse than no result, because it closes an
open question.

\paragraph{Implication.} Intersectional analysis should be the default rather
than an optional deepening, and per-attribute screening should be understood as
the weakest available audit rather than the standard one. This is
\citet{buolamwini2018}'s finding in a lending setting where the marginal audit
does not merely understate the problem --- it returns a clean result. As
\citet{foulds2020} note, the number of subgroups grows combinatorially and
statistical power falls with it, which is precisely why $n$ must be reported
beside every intersectional ratio rather than after it.

\section{Finding 3: the age audit used the wrong statute}
\label{sec:ecoa}

Fairness work on credit routinely audits age at a $40+$ threshold; our pipeline
did. That number comes from the Age Discrimination in Employment Act, which
governs employment. Credit is governed by the Equal Credit Opportunity Act, and
Regulation B \S1002.2(o) defines the protected class as \textbf{62 or older}.
A comment in our own cleaning script asserted that $40+$ ``aligns with US ECOA''.
It does not.

The substitution is not cosmetic. Our banded \texttt{age\_group} attribute pools
young and senior applicants against the middle band. The two move in opposite
directions --- young applicants are approved at a higher rate than the middle
band, senior applicants at a lower one --- so pooling cancels them. The banded
attribute returns $\mathrm{DI} = 0.9791$, severity \textsc{low}: a clean bill of
health. Audited on the statutory cut, the same model on the same rows returns
$\mathbf{0.8897}$, the largest age disparity in the dataset and the lowest of its
four protected attributes (Table~\ref{tab:uc2}).

The flag needed to do this, \texttt{applicant\_age\_above\_62}, was present and
unused in the raw HMDA file. It is published precisely because \S1002.2(o) turns
on it, and it is necessary because the public age bands straddle $62$ --- the
\texttt{55-64} band cannot be split, so no banded attribute can isolate the
protected class. The information was there; the pipeline discarded it in favour
of a threshold imported from the wrong body of law.

One further asymmetry this exposes and cannot resolve: \S1002.6(b)(2) expressly
permits \emph{favouring} an elderly applicant. Disparate impact is symmetric, so
a ratio above $1$ on this attribute is lawful and is not a reverse-disparity
finding, yet the metric reports it identically. The instrument cannot express the
rule it is being used to check, which is the general form of this paper's claim
stated in statute rather than in code.

\section{Finding 4: our own instruments penalised the right answer}
\label{sec:instruments}

The two findings above concern conventions we inherited. This one concerns code
we wrote. Both instruments below were built deliberately, by authors who
understood the domain, to hold the language model to account. Both systematically
punished correct behaviour. We report them unrepaired.

\subsection{The refusal detector measured brevity}
\label{sec:refusal}

Our refusal detector flagged six of nine responses under the
\texttt{constrained} prompt strategy. \emph{None was a refusal.} No flagged
record contained a refusal phrase, a missing field, or an invalid severity value.

The detector treats any narrative field under 40 characters as effectively empty.
The \texttt{constrained} strategy explicitly instructs the model to answer
tersely. So \texttt{how\_to\_fix: "DisparateImpactRemover"} --- a precise,
correct, on-spec answer naming the applicable AIF360 post-processor in 22
characters --- was recorded as a refusal. Compliance with the instruction was the
trigger.

Reporting the raw output would have yielded a $66.7\%$ refusal rate for
\texttt{constrained} and a tidy narrative about constrained prompting inducing
model evasion. That narrative would have been entirely an artifact of our
measuring instrument. \citet{wester2024} study genuine LLM denials, and their
work is the reason we built the detector; our detector does not measure the
phenomenon they describe.

We did not change it. It is part of a frozen scoring harness, and altering it
would invalidate the pre-freeze comparison in \S\ref{sec:interpretive}. Instead
the analysis separates genuine refusals from brevity flags and reports both
columns.

\subsection{The guardrail punished ``collect more data''}

Our guardrail baseline requires every proposed mitigation to name a canonical
algorithm --- reweighing \citep{kamiran2011}, disparate impact removal
\citep{feldman2015}, calibrated equalised odds, and so on. It is a reasonable
gate: it discriminates between a concrete remediation and vague advice, and the
named methods are exactly the ones a practitioner would reach for.

Exactly two of nine attributes failed it. They are
\texttt{foreign\_worker\_original} ($n = 6$ in the privileged group) and
\texttt{loan\_amount\_level} ($3.2\times$ tier imbalance). In both cases the
model proposed acquiring records to a stated floor, Bayesian hierarchical
shrinkage, bootstrapped confidence intervals, and manual review until $n$ is
adequate --- and in both cases that is the correct answer. No post-processor
repairs an $n = 6$ subgroup. Applying one would manufacture false confidence in a
quantity that cannot be estimated.

The gate encodes an assumption: that every fairness problem has an algorithmic
mitigation. Where the problem is data scarcity, the assumption is false, and the
gate penalises the only defensible response. The two failures are not noise ---
they are the two attributes where the assumption does not hold, which is the
strongest form this finding could take.

We first suspected vocabulary matching and normalised the keyword list for
hyphenation, since \texttt{o3} writes ``Calibrated-Equalized-Odds'' and
``Kamiran-Calders Re-weighing''. That fix was necessary and revealed the real
signal rather than removing it: these two survive it. We did not widen the
keyword list further, because doing so would convert a substantive finding into a
passing score.

\subsection{The fabrication detector penalises arithmetic}
\label{sec:fabrication}

Our second detector flags any numeric value in a response that does not appear
in the provided metric context, on the reasonable theory that an auditor should
not invent numbers. Replaying the same packs against two models a generation
apart shows what it actually measures.

\texttt{gpt-5-mini} trips $17$ fabrication flags across $28$ calls;
\texttt{gpt-4o}, on identical input, trips $4$. Read naively, the newer model
fabricates five times as often. Checking each flagged value against the context,
\textbf{$14$ of $21$ ($67\%$) are derivable from it}: $88.88$ is the \texttt{Sex}
disparate impact of $0.8888$ restated as a percentage, $83.87$ is
\texttt{foreign\_worker}'s $0.8387$, and most of the remainder are gaps between
two context values expressed in percentage points. The detector matches raw
float literals, so converting a ratio to a percentage or subtracting two
provided numbers registers as fabrication.

\texttt{gpt-5-mini} does more of this arithmetic because it reasons more
explicitly. The instrument therefore ranks the more capable model as the less
trustworthy one, and it does so specifically because that model shows its work.
Seven values ($19.9$ and $5.3$, each repeated across cycles) were not derivable
by this test and remain candidates for genuine fabrication; their repetition
suggests a systematic derivation we have not identified rather than invention.

This is the same failure as \S\ref{sec:refusal}, found independently in a second
instrument written by the same authors on the same day: a mechanical proxy
scoring a correct behaviour as a fault, in the direction that flatters the
weaker system.

\subsection{Why this belongs in the paper}

Checklists, gates and detectors are the standard response to the problem that
fairness work is hard to verify --- the motivation behind model cards
\citep{mitchell2019}, datasheets \citep{gebru2021} and co-designed checklists
\citep{madaio2020}, and the operational content of internal audit frameworks
\citep{raji2020}. \citet{kroll2021} argues that traceability is what makes
accountability actionable, and we agree.

The cost is that a mechanical proxy stands in for a judgement, and it will be
wrong in the cases where judgement was most needed. Both of our instruments
failed on the hardest inputs: the terse-but-correct answer, and the subgroup too
small to analyse. Neither failure would have been visible from aggregate scores.
Anyone reusing this harness inherits both proxies, which is why we document them
rather than quietly fixing them.

\section{The interpretive layer is contestable}
\label{sec:interpretive}

Having argued that the deterministic apparatus carries commitments, we turn to
the LLM layer built on top of it. The results support a narrow claim: the
interpretive step is real work, it is unreliable in specific measurable ways, and
its aggregate summary statistics should be distrusted.

\subsection{Aggregate agreement is not a meaningful quantity}

Overall severity agreement between \texttt{gpt-4o} and the deterministic labels
is $27/36 = 75.0\%$. That number should not be quoted alone
(Table~\ref{tab:agreement}).

\begin{table}[t]
\caption{Track Q severity agreement. The ordering tracks classifier
discriminativeness, not model capability.}
\label{tab:agreement}
\small
\begin{tabular}{@{}lrr@{}}
\toprule
Use case & Agreed & Rate\\
\midrule
german\_credit & 4/12 & 33.3\%\\
hmda & 11/12 & 91.7\%\\
lending\_club & 12/12 & 100.0\%\\
\midrule
\textbf{Overall} & \textbf{27/36} & \textbf{75.0\%}\\
\bottomrule
\end{tabular}
\end{table}

The ordering is not a competence ranking. UC3's $100\%$ is the easiest possible
case: the classifier predicts non-default for roughly $99.5\%$ of the test set
(four default predictions in 600), so every attribute is a trivial near-parity
\textsc{low} and agreeing costs nothing. UC1's $33.3\%$ is the hardest: three
\textsc{moderate} baselines sitting near threshold boundaries, where a small
interpretive difference flips the label. Aggregate agreement is therefore a
function of how discriminating the underlying classifiers are --- a property of
the benchmark, not of the model. A single headline agreement number for an LLM
fairness auditor is not a meaningful quantity, and we would have reported one had
we not decomposed it.

Errors are asymmetric in the safe direction: seven over-escalations against two
under-escalations. Under-escalation is the consequential failure for an audit
tool; over-escalation costs reviewer time.

\subsection{Prompt phrasing moves the verdict}

Three of nine attribute pairs changed severity label across the four prompt
strategies on \emph{identical} input context.
\texttt{german\_credit/Sex\_original} went
\textsc{high} $\to$ \textsc{high} $\to$ \textsc{moderate} $\to$ \textsc{high};
\texttt{hmda/race} went \textsc{low} $\to$ \textsc{moderate} $\to$
\textsc{moderate} $\to$ \textsc{moderate}. Same model, same numbers, different
wording, different answer. For an audit artifact that may be cited in a
compliance conversation, this is the central reliability result.

\begin{table}[t]
\caption{Track Q by prompt strategy. \texttt{self\_critique} leads on both axes;
\texttt{zero\_shot} carries five of eight hallucination flags.}
\label{tab:strategies}
\small
\begin{tabular}{@{}llrrr@{}}
\toprule
Cycle & Strategy & Mean & Agreement & Halluc.\\
\midrule
1 & \texttt{zero\_shot} & 76.67 & 55.6\% & 5\\
2 & \texttt{chain\_of\_thought} & 80.00 & 77.8\% & 1\\
3 & \texttt{self\_critique} & \textbf{83.89} & \textbf{88.9\%} & 1\\
4 & \texttt{constrained} & 76.56 & 77.8\% & 1\\
\bottomrule
\end{tabular}
\end{table}

\subsection{Grounding suppresses fabrication}

Comparing like-for-like against an earlier run over the same two use cases with
the same model, where the only difference is whether retrieved literature was
embedded in the prompt: fabricated-citation flags fall from $16$ to $2$, and mean
score rises from $69.08$ to $76.67$. This is a design result rather than a model
result, and it is the most actionable finding in the LLM half of the study:
grounding the prompt in real retrieved evidence cuts fabrication roughly
eight-fold.

\subsection{The provider dominates the prompt}

Running the identical harness against \texttt{gemini-2.5-flash} produced refusals
in $19$ of $24$ cycles, a mean score of $19.79$, and $14$ zero-score cycles. No
prompt strategy compensates for that. Any benchmark of this shape measures
provider capability at least as much as it measures method --- a caveat that
applies to our own numbers, and one reason we release the frozen prompts rather
than only the scores.

\subsection{What the narrative form added}

Track H audits passed the citation check at $27/27$ against the retrieved
evidence pool, with weighted rubric scores of $4.8$, $4.7$ and $5.0$ out of $5$
(cross-model: \texttt{o3} generated, \texttt{gpt-4o} graded, reported as
secondary evidence because a model grading a model is weaker than a mechanical
gate \citep{zheng2023}).

Two behaviours are worth naming because the deterministic template cannot produce
them. On \texttt{foreign\_worker\_original} the audit independently reproduced the
$n = 6$ caveat and quantified it: one label change would swing DI to $0.78$,
therefore statistical power is inadequate. On UC3 it identified the
near-degenerate classifier unprompted and declined the flattering reading,
converting a null result into a conditional prediction with a mechanism and a
supporting citation --- that near-parity here reflects a classifier approving
almost everything, and that latent disparities may surface once realistic credit
thresholds are enforced.

The deterministic pipeline establishes what is true. The narrative layer
establishes what it means for a decision. They are not competing
implementations, and asking which is more accurate is a category error.

\section{Discussion}

\paragraph{Audit infrastructure needs its own audit.} Our three findings share a
structure: a component that is correct in its own terms produces a wrong or
misleading conclusion at the level of the audit. Metric libraries, severity
thresholds, refusal detectors and quality gates are all measurement instruments,
and measurement instruments require validation against cases where the right
answer is known independently. We had that only by accident --- the $n = 6$
subgroup and the 22-character correct answer were adversarial test cases we did
not design. Audit tooling should ship with them deliberately.

\paragraph{Conventions should become parameters.} The orientation failure
(\S\ref{sec:orientation}) is repairable by design: make the favorable outcome a
required argument. Generalising, any place where a fairness pipeline infers a
normative category from a syntactic feature --- a label value, a column name, a
character count --- is a place where the pipeline will silently encode someone's
assumption. Making these explicit costs one parameter and converts a class of
silent inversions into visible errors.

\paragraph{Which attributes deserve protection is deferred, not answered.} UC3
audits \texttt{income\_level} and \texttt{loan\_amount\_level} beside
\texttt{gender} with identical machinery and identical thresholds. Near-parity on
gender and near-parity on income level are not the same kind of result: one
concerns a protected class, the other an economic characteristic that lending law
generally permits pricing on. The pipeline cannot tell them apart, and by
treating them uniformly it silently defers a normative question to whoever
configured the attribute list. This is the point at which the stakeholder
literature \citep{luo2025,luo2026,deshpande2022} becomes not merely
complementary but necessary: the apparatus has no resources for the question.

\paragraph{Thresholds are screens, not verdicts.} We report DI against the
four-fifths screen throughout, and both \S\ref{sec:intersectional} readings turn
on refusing to treat it as a gate. A failed screen with EOD $\approx 0$ means
label bias; a passed screen at $n = 6$ means nothing at all.
\citet{barocas2016} and \citet{wachter2021} establish why the legal question does
not reduce to the ratio, and our results show the technical question does not
either.

\paragraph{For practitioners.} Three concrete moves follow: state the favorable
outcome explicitly and test it; compute intersectional cells by default and print
$n$ beside every ratio; and treat each mechanical gate in the audit stack as a
hypothesis to be falsified, starting with the inputs where the correct answer is
unusual.

\section{Limitations}

\textbf{Subgroup sizes.} \texttt{foreign\_worker\_original} has $n = 6$ in the
privileged group and several HMDA race categories are severely sparse.
Intersectional cells are reported with $n$ and excluded from conclusions below $n = 15$ --- a floor whose cost we quantify in \S\ref{sec:floor}.

\textbf{UC3 is barely informative as a fairness case.} A classifier predicting
one class $\sim99.5\%$ of the time cannot exhibit measurable disparity. UC3
exercises the pipeline; it does not audit a working model. Its value here is the
orientation artifact and the proxy-versus-class contrast.

\textbf{LLM nondeterminism.} We pinned neither temperature nor seed. The
four-cycle design varies prompt strategy, not repeated samples of one strategy,
so cycle-to-cycle differences confound strategy with sampling noise. The
severity-flip result in \S\ref{sec:interpretive} should be read as ``phrasing or
sampling'' rather than phrasing alone.

\textbf{API versioning.} \texttt{gpt-4o} and \texttt{o3} are moving targets.
Frozen prompt packs make the inputs reproducible; the outputs are not.

\textbf{Metric grounding.} All metrics are associational. Root-cause mapping is
rule-based inference over metric values and feature correlations, not causal
identification. We make no causal claim.

\textbf{Cost accounting.} \texttt{o3} is priced at the higher of two conflicting
internal rate tables, so reported spend is an upper bound rather than a billed
amount.

\textbf{No legal claim.} ECOA, Regulation B and the four-fifths rule appear as
decision context. A disparate impact ratio below $0.8$ is a screening signal, not
a finding of unlawful discrimination \citep{barocas2016,wachter2021}.

\textbf{Scope.} HMDA is Georgia only; German Credit is 1990s German data used as
a benchmark rather than a current population. Three datasets, one primary model,
and a single institution's judgement about what counts as a correct mitigation.

\textbf{A known artifact defect.} Our consolidated metric CSV carries two
disjoint column schemas because per-dataset files are concatenated without
harmonisation. The per-dataset files are authoritative. We note it because it is
exactly the class of infrastructure defect this paper is about.

\section{Conclusion}

We set out to benchmark whether a language model can interpret fairness metrics
and found that the harder problem was upstream. A disparate impact ratio of $0.0$
that should have been $0.9931$; three passing marginals concealing an
intersectional ratio of $0.6087$; a refusal detector that measured brevity; a
guardrail that penalised the only correct answer to a six-record subgroup. None
of these was a bug in the ordinary sense, and none was user error. Each was a
defensible design decision whose normative content became visible only when it
punished a correct answer.

If fairness is defined relative to a human decision context
\citep{lee2020}, then the apparatus that measures fairness inherits that
dependence. Our instruments were built by people who knew the domain and still
encoded assumptions that inverted conclusions. We expect this to be the common
case rather than the exception, and we release the frozen artifacts so that the
claim can be checked against our own instruments rather than taken on our word.

\begin{acks}
We thank the maintainers of AIF360 and Fairlearn. Retrieved literature was
obtained via the Semantic Scholar API.
\end{acks}

\bibliographystyle{ACM-Reference-Format}
\bibliography{refs}

\end{document}
